% ============================================================
\chapter{Syst\`emes Dynamiques Discrets}
\label{ch:discrets}
% ============================================================

\section{Introduction}

En 1976, le biologiste Robert May publie un article dans \emph{Nature} montrant que l'application logistique $x_{n+1} = r x_n(1 - x_n)$, un mod\`ele \'el\'ementaire de dynamique de population, produit un comportement d'une richesse stupéfiante~: points fixes stables, cycles de p\'eriode $2$, $4$, $8$\ldots{}, puis chaos. Mitchell Feigenbaum d\'ecouvre bient\^ot que les rapports entre les valeurs du param\`etre o\`u se produisent les doublements de p\'eriode convergent vers une constante universelle, $\delta \approx 4{,}669$, ind\'ependante de la famille d'applications consid\'er\'ee. Cette \emph{universalit\'e} relie la dynamique discr\`ete \`a la physique statistique. Ce chapitre explore les syst\`emes d\'efinis par l'it\'eration $x_{n+1} = f(x_n)$, le cadre le plus simple pour \'etudier le chaos~: cascades de doublements, dynamique symbolique et entropie topologique.

\begin{intuition}
M\^eme une application aussi simple que $f(x) = rx(1-x)$ engendre, lorsque
$r$ augmente, toute la complexit\'e du chaos : bifurcations, auto-similarit\'e,
universalit\'e. Les syst\`emes discrets sont le laboratoire id\'eal pour
comprendre les m\'ecanismes fondamentaux.
\end{intuition}

\section{L'application logistique}

\begin{definition}[Famille logistique]
\index{application logistique}
La \textbf{famille logistique} est $f_r : [0,1] \to [0,1]$ d\'efinie par
$f_r(x) = rx(1-x)$, $r \in [0, 4]$.
\end{definition}

\begin{proposition}[Points fixes et stabilit\'e]
\begin{enumerate}
    \item Le point fixe $x^* = 0$ existe pour tout $r$ ; il est stable si $r < 1$.
    \item Le point fixe $x^* = 1 - 1/r$ existe pour $r > 1$ ; il est stable si
    $1 < r < 3$ (car $f_r'(x^*) = 2 - r$).
    \item \`A $r = 3$, une bifurcation de doublement de p\'eriode cr\'ee un
    cycle de p\'eriode 2.
\end{enumerate}
\end{proposition}

\begin{theoreme}[Cascade de doublements de p\'eriode]
\index{cascade de Feigenbaum}
Il existe une suite croissante $r_1 < r_2 < r_3 < \cdots$ avec
$r_1 = 3$ telle qu'en $r = r_n$, le cycle de p\'eriode $2^{n-1}$ perd sa
stabilit\'e et un cycle stable de p\'eriode $2^n$ appara\^it par bifurcation
flip. La suite $r_n$ converge vers $r_\infty \approx 3{,}5699\ldots$
\end{theoreme}

\begin{theoreme}[Universalit\'e de Feigenbaum]
\index{constante de Feigenbaum}
Le rapport $\delta_n = \frac{r_n - r_{n-1}}{r_{n+1} - r_n}$ converge vers
la constante universelle de Feigenbaum
\[
    \delta_F = \lim_{n \to \infty} \delta_n = 4{,}669\,201\,609\ldots
\]
Cette constante est la m\^eme pour toute famille unimodale \`a maximum
quadratique. Le facteur d'\'echelle spatial converge vers
$\alpha_F = -2{,}502\,907\,875\ldots$
\end{theoreme}

\begin{formules}[title=\textbf{R\'egimes de l'application logistique}]
\begin{align*}
    r \in [0, 1] &: \text{convergence vers } 0. \\
    r \in (1, 3) &: \text{convergence vers } x^* = 1 - 1/r. \\
    r \in (3, 1+\sqrt{6}) &: \text{cycle stable de p\'eriode } 2. \\
    r = r_\infty \approx 3{,}5699 &: \text{accumulation des doublements.} \\
    r \in (r_\infty, 4) &: \text{m\'elange de chaos et fen\^etres p\'eriodiques.} \\
    r = 4 &: \text{chaos d\'evelopp\'e sur } [0, 1].
\end{align*}
\end{formules}

\section{Le th\'eor\`eme de Sharkovskii}

\begin{definition}[Ordre de Sharkovskii]
\index{ordre de Sharkovskii}
L'\textbf{ordre de Sharkovskii} sur $\N^*$ est :
\begin{multline*}
    3 \triangleright 5 \triangleright 7 \triangleright 9 \triangleright \cdots
    \triangleright 2 \cdot 3 \triangleright 2 \cdot 5 \triangleright 2 \cdot 7
    \triangleright \cdots \\
    \triangleright 4 \cdot 3 \triangleright 4 \cdot 5 \triangleright \cdots
    \triangleright \cdots \triangleright 2^3 \triangleright 2^2
    \triangleright 2 \triangleright 1.
\end{multline*}
\end{definition}

\begin{theoreme}[Sharkovskii, 1964]
\index{th\'eor\`eme de Sharkovskii}
Soit $f : \R \to \R$ continue. Si $f$ a un point p\'eriodique de p\'eriode $m$,
alors $f$ a des points p\'eriodiques de toute p\'eriode $n$ telle que
$m \triangleright n$ dans l'ordre de Sharkovskii. De plus, pour tout $m$,
il existe une application continue ayant un point de p\'eriode $m$ mais aucun
point de p\'eriode $n$ avec $n \triangleright m$.
\end{theoreme}

\begin{corollaire}[P\'eriode 3 implique toutes les p\'eriodes]
Si $f : \R \to \R$ continue a un point de p\'eriode 3, alors $f$ a des points
p\'eriodiques de toute p\'eriode $n \in \N^*$.
\end{corollaire}

\begin{remarque}
L'ordre de Sharkovskii ne s'applique qu'en dimension 1. En dimension 2,
un diff\'eomorphisme peut avoir un point de p\'eriode 3 sans avoir de point
fixe (rotation d'angle $2\pi/3$).
\end{remarque}

\begin{lemme}[Lemme du graphe]
\label{lem:graphe}
Soient $I, J$ des intervalles ferm\'es et $f$ continue. Si $f(I) \supset J$,
alors il existe un sous-intervalle $K \subset I$ tel que $f(K) = J$.
En particulier, si $f(I) \supset I$, alors $f$ a un point fixe dans $I$.
\end{lemme}

\section{Dynamique symbolique}

\begin{definition}[Espace des suites]
\index{dynamique symbolique}
Soit $\mathcal{A} = \{0, 1, \ldots, k-1\}$ un alphabet de $k$ symboles.
L'\textbf{espace des suites} est
$\Sigma_k = \mathcal{A}^\N = \{(s_0, s_1, s_2, \ldots) : s_i \in \mathcal{A}\}$,
muni de la m\'etrique $d(s, t) = \sum_{i=0}^{\infty} \frac{\abs{s_i - t_i}}{k^i}$.
\end{definition}

\begin{definition}[Application de d\'ecalage (shift)]
\index{shift}
Le \textbf{shift} $\sigma : \Sigma_k \to \Sigma_k$ est d\'efini par
$\sigma(s_0, s_1, s_2, \ldots) = (s_1, s_2, s_3, \ldots)$.
\end{definition}

\begin{theoreme}[Propri\'et\'es du shift]
Le shift $\sigma$ sur $\Sigma_k$ satisfait :
\begin{enumerate}
    \item $\sigma$ est continue et surjective (mais non injective).
    \item Les points p\'eriodiques de p\'eriode $n$ sont exactement les suites
    $n$-p\'eriodiques ; il y en a $k^n$.
    \item Les points p\'eriodiques sont denses dans $\Sigma_k$.
    \item $\sigma$ est topologiquement transitive.
    \item $\sigma$ est chaotique au sens de Devaney.
\end{enumerate}
\end{theoreme}

\begin{definition}[Sous-shift de type fini]
\index{sous-shift de type fini}
\'Etant donn\'ee une matrice de transition $A \in \{0,1\}^{k \times k}$, le
\textbf{sous-shift de type fini} est
$\Sigma_A = \{s \in \Sigma_k : A_{s_i, s_{i+1}} = 1 \text{ pour tout } i\}$,
et $\sigma_A = \sigma|_{\Sigma_A}$.
\end{definition}

\begin{proposition}[Codage des it\'er\'ees]
Soit $f : [0,1] \to [0,1]$ unimodale avec point critique $c$. L'\textbf{itin\'eraire}
de $x$ est la suite $s(x) = (s_0, s_1, \ldots)$ o\`u $s_n = 0$ si $f^n(x) < c$
et $s_n = 1$ si $f^n(x) > c$. L'application $x \mapsto s(x)$ semi-conjugue
$f$ au shift sur un sous-espace de $\Sigma_2$.
\end{proposition}

\section{Entropie topologique}

\begin{definition}[Entropie topologique]
\index{entropie topologique}
Soit $f : X \to X$ une application continue sur un espace m\'etrique compact.
L'\textbf{entropie topologique} est
\[
    h_{\mathrm{top}}(f) = \lim_{n \to \infty} \frac{1}{n} \ln s(n, \epsilon, f),
\]
o\`u $s(n, \epsilon, f)$ est le nombre maximal de points
$(n, \epsilon)$-s\'epar\'es : pour tous $x \neq y$ dans l'ensemble,
$\max_{0 \leq k < n} d(f^k(x), f^k(y)) > \epsilon$.
La limite est ind\'ependante de $\epsilon > 0$ suffisamment petit.
\end{definition}

\begin{theoreme}[Entropie du shift]
L'entropie topologique du shift $\sigma$ sur $\Sigma_k$ est
$h_{\mathrm{top}}(\sigma) = \ln k$.
Pour le sous-shift de type fini $\sigma_A$, $h_{\mathrm{top}}(\sigma_A) = \ln \rho(A)$,
o\`u $\rho(A)$ est le rayon spectral de $A$.
\end{theoreme}

\begin{proposition}[Entropie de l'application logistique]
Pour $f_4(x) = 4x(1-x)$, l'entropie topologique est $h_{\mathrm{top}}(f_4) = \ln 2$.
Plus g\'en\'eralement, pour l'application de la tente de pente $s$ :
$h_{\mathrm{top}} = \max(0, \ln s)$.
\end{proposition}

\begin{theoreme}[In\'egalit\'e variationnelle]
\index{principe variationnel}
Pour toute application continue $f$ sur un espace m\'etrique compact :
\[
    h_{\mathrm{top}}(f) = \sup_{\mu \in \mathcal{M}_f} h_\mu(f),
\]
o\`u le supremum porte sur toutes les mesures de probabilit\'e $f$-invariantes
et $h_\mu(f)$ est l'entropie m\'etrique.
\end{theoreme}

\begin{attention}[title=\textbf{Entropie et complexit\'e}]
$h_{\mathrm{top}} > 0$ implique que le nombre d'orbites distinguables cro\^it
exponentiellement avec le temps. C'est un indicateur de chaos plus robuste que
l'exposant de Lyapunov maximal, car il est un invariant topologique (ne d\'epend
pas de la m\'etrique ni de la mesure).
\end{attention}

\section{Exemples avanc\'es}

\begin{exemple}[Application en dents de scie (Bernoulli)]
L'application $B_k(x) = kx \pmod{1}$ sur $[0, 1)$ est exactement conjugu\'ee
au shift sur $\Sigma_k$ via le d\'eveloppement en base $k$. Son entropie est
$\ln k$ et son exposant de Lyapunov est $\ln k$.
\end{exemple}

\begin{exemple}[Ensemble de Julia quadratique]
Pour $f_c(z) = z^2 + c$ avec $c \in \C$, l'ensemble de Julia
$J_c = \partial\{z : f_c^n(z) \not\to \infty\}$ est soit connexe (si $c$ est
dans l'ensemble de Mandelbrot) soit totalement discontinu (un ensemble de
Cantor). La dynamique sur $J_c$ est riche en ph\'enom\`enes chaotiques.
\end{exemple}

\section{Exercices}

\begin{exercice}[Bifurcations logistiques]
Pour l'application logistique $f_r(x) = rx(1-x)$, trouver analytiquement la
valeur de $r$ \`a laquelle le cycle de p\'eriode 2 appara\^it et les valeurs
des points du cycle.
\end{exercice}

\begin{exercice}[Sharkovskii]
Construire une application continue $f : [0,1] \to [0,1]$ ayant un point de
p\'eriode 6 mais aucun point de p\'eriode 3. Montrer qu'elle a n\'ecessairement
des points de p\'eriode 2.
\end{exercice}

\begin{exercice}[Dynamique symbolique]
Montrer que le shift sur $\Sigma_2$ est topologiquement m\'elangeant : pour
tous ouverts $U, V$, il existe $N$ tel que $\sigma^n(U) \cap V \neq \emptyset$
pour tout $n \geq N$.
\end{exercice}

\begin{exercice}[Entropie]
Calculer l'entropie topologique du sous-shift de type fini d\'efini par la
matrice $A = \begin{pmatrix} 1 & 1 \\ 1 & 0 \end{pmatrix}$ (shift de Fibonacci).
Montrer que $h_{\mathrm{top}} = \ln\varphi$, o\`u $\varphi = (1+\sqrt{5})/2$.
\end{exercice}

\begin{exercice}[Diagramme de bifurcation]
\'Ecrire un programme tra\c{c}ant le diagramme de bifurcation de l'application
logistique pour $r \in [2{,}5, \, 4]$. Identifier visuellement la cascade de
doublements et les fen\^etres p\'eriodiques (p\'eriode 3, 5, 6, etc.).
\end{exercice}

\begin{exercice}[Itin\'eraire de kneading]
D\'eterminer l'itin\'eraire de kneading du point critique de $f_r$ pour
$r = 3{,}8$ et $r = 4$. En d\'eduire quelles p\'eriodes sont forc\'ees.
\end{exercice}
